\subsection{A.1 Vanilla RNN με Attention}
 
\noindent \textbf{Πηγή:} \texttt{Vanilla\_RNN\_Temporal\_Protein.ipynb}
 
 
\subsubsection*{Temporal Protein Classification with Vanilla RNN}
 
\subsubsection*{ΜΕΡΟΣ 1: SETUP \& IMPORTS}
 
\begin{lstlisting}
# Imports
import numpy as np
import pandas as pd
import torch
import torch.nn as nn
import torch.nn.functional as F
from torch.utils.data import Dataset, DataLoader
from sklearn.model_selection import StratifiedKFold
from sklearn.preprocessing import StandardScaler, LabelEncoder
from sklearn.feature_selection import SelectKBest, f_classif
from sklearn.impute import SimpleImputer
from sklearn.metrics import classification_report, confusion_matrix, accuracy_score
import matplotlib.pyplot as plt
import seaborn as sns
import warnings
warnings.filterwarnings('ignore')
 
# Device selection
device = torch.device('cuda' if torch.cuda.is_available() else 'cpu')
print(f"="*50)
print(f"Using device: {device}")
if torch.cuda.is_available():
    print(f"GPU: {torch.cuda.get_device_name(0)}")
print(f"="*50)
 
# Reproducibility
def set_seed(seed=42):
    """Set seed for reproducibility"""
    np.random.seed(seed)
    torch.manual_seed(seed)
    if torch.cuda.is_available():
        torch.cuda.manual_seed_all(seed)
    torch.backends.cudnn.deterministic = True
    torch.backends.cudnn.benchmark = False
 
set_seed(42)
print("\n✓ Setup complete!")
\end{lstlisting}
 
\subsubsection*{ΜΕΡΟΣ 2: CONFIGURATION}
 
\begin{lstlisting}
class Config:
    """
    Configuration class - Όλες οι hyperparameters εδώ
    """
    # === Random Seed ===
    RANDOM_SEED = 42
 
    # === Data ===
    N_FOLDS = 5                    # Cross-validation folds
    N_TOP_FEATURES = 50            # Πόσες πρωτεΐνες να κρατήσουμε
 
    # === Model Architecture ===
    # CNN
    CNN_FILTERS = 32               # Αριθμός filters στο CNN
    CNN_KERNEL_SIZE = 5            # Kernel size
 
    # RNN
    RNN_HIDDEN_SIZE = 32           # Hidden size του RNN
    RNN_NUM_LAYERS = 1             # Αριθμός RNN layers
    RNN_BIDIRECTIONAL = True       # Bidirectional RNN
 
    # Attention
    ATTENTION_DIM = 32             # Attention dimension
 
    # Regularization
    DROPOUT = 0.5                  # Dropout rate
 
    # === Training ===
    BATCH_SIZE = 8                 # Batch size
    LEARNING_RATE = 1e-4           # Learning rate (μικρό για stability)
    WEIGHT_DECAY = 1e-3            # L2 regularization
    EPOCHS = 100                   # Max epochs
    PATIENCE = 15                  # Early stopping patience
 
    # === Device ===
    DEVICE = device
 
# Print configuration
print("Configuration:")
print(f"  Features: {Config.N_TOP_FEATURES}")
print(f"  RNN Hidden: {Config.RNN_HIDDEN_SIZE}")
print(f"  Bidirectional: {Config.RNN_BIDIRECTIONAL}")
print(f"  Dropout: {Config.DROPOUT}")
print(f"  Learning Rate: {Config.LEARNING_RATE}")
print(f"  Cross-validation Folds: {Config.N_FOLDS}")
print("\n✓ Configuration loaded!")
\end{lstlisting}
 
\subsubsection*{ΜΕΡΟΣ 3: DATA LOADING}
 
\begin{lstlisting}
# ╔══════════════════════════════════════════════════════════════╗
# ║  ΑΛΛΑΞΕ ΑΥΤΑ ΤΑ PATHS ΣΤΑ ΔΙΚΑ ΣΟΥ ΑΡΧΕΙΑ                   ║
# ╚══════════════════════════════════════════════════════════════╝
 
SAMPLES_PATH = '/content/DIOPTRA_Samples_list_Oncology.xlsx'
PROTEINS_PATH = '/content/DIOPTRA_Protein_results_Oncology.xlsx'
\end{lstlisting}
 
\begin{lstlisting}
def load_data(samples_path, proteins_path, include_t00=True):
    """
    Φορτώνει και προετοιμάζει τα δεδομένα.
 
    Parameters:
    -----------
    samples_path : str
        Path στο αρχείο με τα sample metadata
    proteins_path : str
        Path στο αρχείο με τις τιμές πρωτεϊνών
    include_t00 : bool
        Αν θα συμπεριλάβουμε τα single-timepoint samples (T00)
 
    Returns:
    --------
    X : np.ndarray
        Shape (n_samples, 2, n_proteins)
    y : np.ndarray
        Labels
    protein_names : list
        Ονόματα πρωτεϊνών
    label_encoder : LabelEncoder
        Fitted encoder
    """
    print("Loading data...")
    print("="*50)
 
    # Load Excel files
    samples = pd.read_excel(samples_path)
    proteins = pd.read_excel(proteins_path)
 
    print(f"Samples file: {len(samples)} rows")
    print(f"Proteins file: {len(proteins)} rows, {len(proteins.columns)-1} proteins")
 
    # Merge data
    data = samples.merge(proteins, on='SampleID')
 
    # Get protein column names
    protein_cols = [c for c in proteins.columns if c != 'SampleID']
 
    # === Create paired samples (T0 → T1) ===
    print("\nCreating paired samples...")
 
    t0_data = data[data['Timepoint'] == 'T0'].copy()
    t1_data = data[data['Timepoint'] == 'T1'].copy()
 
    # Patient ID extraction
    t0_data['PatientID'] = t0_data['SampleID']
    t1_data['PatientID'] = t1_data['SampleID'].str.replace('-F', '', regex=False)
 
    # Merge T0 and T1
    paired = t0_data.merge(
        t1_data[['PatientID'] + protein_cols],
        on='PatientID',
        suffixes=('_T0', '_T1')
    )
    print(f"  Paired samples (T0→T1): {len(paired)}")
 
    # === Add T00 samples (single timepoint) ===
    if include_t00:
        t00_data = data[data['Timepoint'] == 'T00'].copy()
 
        # Duplicate as T0 and T1 (no temporal change)
        for col in protein_cols:
            t00_data[f'{col}_T0'] = t00_data[col]
            t00_data[f'{col}_T1'] = t00_data[col]
        t00_data['PatientID'] = t00_data['SampleID']
 
        # Combine
        cols_needed = ['PatientID', 'Group'] + [f'{c}_T0' for c in protein_cols] + [f'{c}_T1' for c in protein_cols]
        combined = pd.concat([paired[cols_needed], t00_data[cols_needed]], ignore_index=True)
        print(f"  Single timepoint (T00): {len(t00_data)}")
        print(f"  Total samples: {len(combined)}")
    else:
        combined = paired
 
    # === Remove NAA class (too few samples) ===
    combined = combined[combined['Group'] != 'NAA']
    print(f"  After removing NAA: {len(combined)}")
 
    # === Create numpy arrays ===
    X_T0 = combined[[f'{c}_T0' for c in protein_cols]].values
    X_T1 = combined[[f'{c}_T1' for c in protein_cols]].values
    X = np.stack([X_T0, X_T1], axis=1)  # Shape: (n_samples, 2, n_proteins)
 
    # === Encode labels ===
    label_encoder = LabelEncoder()
    y = label_encoder.fit_transform(combined['Group'].values)
 
    # === Print class distribution ===
    print(f"\nClass distribution:")
    for i, cls in enumerate(label_encoder.classes_):
        count = (y == i).sum()
        pct = count / len(y) * 100
        print(f"  {cls}: {count} samples ({pct:.1f}%)")
 
    print(f"\nFinal shapes:")
    print(f"  X: {X.shape} (samples, timepoints, proteins)")
    print(f"  y: {y.shape}")
 
    return X, y, protein_cols, label_encoder
\end{lstlisting}
 
\begin{lstlisting}
# Load the data
X_raw, y, protein_names, label_encoder = load_data(SAMPLES_PATH, PROTEINS_PATH)
class_names = list(label_encoder.classes_)
 
print("\n" + "="*50)
print("✓ Data loaded successfully!")
print(f"  Classes: {class_names}")
\end{lstlisting}
 
\subsubsection*{ΜΕΡΟΣ 4: DATA PREPROCESSING}
 
\subsubsection*{4.1 Handle NaN Values}
 
\begin{lstlisting}
def handle_nan_values(X, protein_names):
    """
    Χειρίζεται τα NaN values:
    1. Αφαιρεί πρωτεΐνες που έχουν μόνο NaN
    2. Κάνει imputation με median στα υπόλοιπα NaN
 
    Parameters:
    -----------
    X : np.ndarray
        Shape (n_samples, n_timepoints, n_proteins)
    protein_names : list
        Ονόματα πρωτεϊνών
 
    Returns:
    --------
    X_clean : np.ndarray
        Cleaned data
    clean_protein_names : list
        Updated protein names
    """
    print("Handling NaN values...")
    print("="*50)
 
    n_samples, n_timepoints, n_proteins = X.shape
    total_nan = np.isnan(X).sum()
    print(f"  Total NaN values: {total_nan}")
    print(f"  Original shape: {X.shape}")
 
    if total_nan == 0:
        print("  No NaN values found!")
        return X, protein_names
 
    # Step 1: Find proteins with ALL NaN values
    valid_protein_mask = ~np.all(np.isnan(X), axis=(0, 1))
    n_all_nan = (~valid_protein_mask).sum()
    print(f"  Proteins with ALL NaN (removed): {n_all_nan}")
 
    # Remove all-NaN proteins
    X = X[:, :, valid_protein_mask]
    protein_names = [p for p, valid in zip(protein_names, valid_protein_mask) if valid]
 
    print(f"  Shape after removing all-NaN proteins: {X.shape}")
 
    # Step 2: Impute remaining NaN with median
    n_samples, n_timepoints, n_proteins = X.shape
    X_imputed = np.zeros_like(X)
 
    imputer = SimpleImputer(strategy='median')
 
    for t in range(n_timepoints):
        X_imputed[:, t, :] = imputer.fit_transform(X[:, t, :])
 
    remaining_nan = np.isnan(X_imputed).sum()
    print(f"  NaN after imputation: {remaining_nan}")
    print(f"  Final shape: {X_imputed.shape}")
 
    return X_imputed, protein_names
\end{lstlisting}
 
\begin{lstlisting}
# Apply NaN handling
X_clean, protein_names = handle_nan_values(X_raw, protein_names)
print("\n✓ NaN handling complete!")
\end{lstlisting}
 
\subsubsection*{4.2 Feature Selection}
 
\begin{lstlisting}
def select_top_features(X, y, protein_names, n_features=50):
    """
    Επιλέγει τις top-K πρωτεΐνες με βάση ANOVA F-score.
 
    Γιατί feature selection:
    - 368 features με 62 samples = overfitting
    - Κρατάμε μόνο τις πιο discriminative
 
    Parameters:
    -----------
    X : np.ndarray
        Shape (n_samples, n_timepoints, n_proteins)
    y : np.ndarray
        Labels
    protein_names : list
        Ονόματα πρωτεϊνών
    n_features : int
        Πόσα features να κρατήσουμε
 
    Returns:
    --------
    X_selected : np.ndarray
        Selected features
    selected_proteins : list
        Names of selected proteins
    selector : SelectKBest
        Fitted selector
    """
    print("Feature Selection...")
    print("="*50)
 
    # Use T0 for feature selection
    X_t0 = X[:, 0, :]
 
    # Adjust n_features if needed
    n_available = X_t0.shape[1]
    n_features = min(n_features, n_available)
    print(f"  Available proteins: {n_available}")
    print(f"  Selecting top: {n_features}")
 
    # ANOVA F-test
    selector = SelectKBest(f_classif, k=n_features)
    selector.fit(X_t0, y)
 
    # Get selected indices
    selected_mask = selector.get_support()
    selected_indices = np.where(selected_mask)[0]
 
    # Apply to both timepoints
    X_selected = X[:, :, selected_indices]
 
    # Get selected protein names
    selected_proteins = [protein_names[i] for i in selected_indices]
 
    # Get scores for ranking
    scores = selector.scores_[selected_indices]
    protein_scores = sorted(zip(selected_proteins, scores), key=lambda x: x[1], reverse=True)
 
    print(f"\n  Top 10 discriminative proteins:")
    for i, (prot, score) in enumerate(protein_scores[:10]):
        print(f"    {i+1}. {prot:<15} (F-score: {score:.2f})")
 
    print(f"\n  Final shape: {X_selected.shape}")
 
    return X_selected, selected_proteins, selector, selected_indices
\end{lstlisting}
 
\begin{lstlisting}
# Apply feature selection
X_selected, selected_proteins, feature_selector, selected_indices = select_top_features(
    X_clean, y, protein_names, n_features=Config.N_TOP_FEATURES
)
print("\n✓ Feature selection complete!")
\end{lstlisting}
 
\subsubsection*{4.3 Normalization}
 
\begin{lstlisting}
def normalize_data(X):
    """
    Κανονικοποιεί τα δεδομένα (mean=0, std=1).
 
    Parameters:
    -----------
    X : np.ndarray
        Shape (n_samples, n_timepoints, n_features)
 
    Returns:
    --------
    X_normalized : np.ndarray
        Normalized data
    scaler : StandardScaler
        Fitted scaler
    """
    print("Normalizing data...")
    print("="*50)
 
    n_samples, n_timepoints, n_features = X.shape
 
    # Flatten to 2D for normalization
    X_flat = X.reshape(-1, n_features)
 
    # Normalize
    scaler = StandardScaler()
    X_flat_scaled = scaler.fit_transform(X_flat)
 
    # Clip extreme values (για numerical stability)
    X_flat_scaled = np.clip(X_flat_scaled, -5, 5)
 
    # Reshape back
    X_normalized = X_flat_scaled.reshape(n_samples, n_timepoints, n_features)
 
    print(f"  Original range: [{X.min():.2f}, {X.max():.2f}]")
    print(f"  Normalized range: [{X_normalized.min():.2f}, {X_normalized.max():.2f}]")
    print(f"  Mean: {X_normalized.mean():.4f}")
    print(f"  Std: {X_normalized.std():.4f}")
 
    return X_normalized, scaler
\end{lstlisting}
 
\begin{lstlisting}
# Apply normalization
X, scaler = normalize_data(X_selected)
print("\n✓ Normalization complete!")
print(f"\nFinal data shape: {X.shape}")
print(f"Any NaN: {np.isnan(X).any()}")
\end{lstlisting}
 
\subsubsection*{ΜΕΡΟΣ 5: DATASET CLASS}
 
\begin{lstlisting}
class ProteinDataset(Dataset):
    """
    PyTorch Dataset για temporal protein data.
 
    Modes:
    - 'temporal': [T0, T1] - πλήρης χρονική πληροφορία
    - 'static': [T0] μόνο - baseline
    - 'delta': [T1-T0] - μόνο η αλλαγή
    """
 
    def __init__(self, X, y, mode='temporal'):
        """
        Parameters:
        -----------
        X : np.ndarray
            Shape (n_samples, 2, n_features)
        y : np.ndarray
            Labels
        mode : str
            'temporal', 'static', or 'delta'
        """
        self.mode = mode
        self.y = torch.LongTensor(y)
 
        if mode == 'temporal':
            # Both timepoints: (n_samples, 2, n_features)
            self.X = torch.FloatTensor(X)
        elif mode == 'static':
            # Only T0: (n_samples, 1, n_features)
            self.X = torch.FloatTensor(X[:, 0:1, :])
        elif mode == 'delta':
            # T1 - T0: (n_samples, 1, n_features)
            delta = X[:, 1, :] - X[:, 0, :]
            self.X = torch.FloatTensor(delta[:, np.newaxis, :])
        else:
            raise ValueError(f"Unknown mode: {mode}. Use 'temporal', 'static', or 'delta'")
 
    def __len__(self):
        return len(self.y)
 
    def __getitem__(self, idx):
        return self.X[idx], self.y[idx]
 
 
# Test dataset
print("Testing Dataset class...")
for mode in ['temporal', 'static', 'delta']:
    ds = ProteinDataset(X, y, mode=mode)
    x_sample, y_sample = ds[0]
    print(f"  {mode:10s}: X shape = {x_sample.shape}")
 
print("\n✓ Dataset class ready!")
\end{lstlisting}
 
\subsubsection*{ΜΕΡΟΣ 6: VANILLA RNN MODEL}
 
\subsubsection*{6.1 Vanilla RNN Cell (Educational Implementation)}
 
\begin{lstlisting}
class VanillaRNNCell(nn.Module):
    """
    Vanilla RNN Cell - Educational Implementation
 
    Εξίσωση:
    h_t = tanh(W_ih · x_t + W_hh · h_{t-1} + b)
 
    Όπου:
    - x_t: input στο χρόνο t
    - h_{t-1}: hidden state από το προηγούμενο βήμα
    - W_ih: weights για το input
    - W_hh: weights για το hidden state
    - b: bias
    - tanh: activation function
    """
 
    def __init__(self, input_size, hidden_size):
        super().__init__()
 
        self.input_size = input_size
        self.hidden_size = hidden_size
 
        # Weights
        self.W_ih = nn.Linear(input_size, hidden_size)   # Input to hidden
        self.W_hh = nn.Linear(hidden_size, hidden_size)  # Hidden to hidden
 
    def forward(self, x, h_prev):
        """
        Forward pass για ένα timestep.
 
        Parameters:
        -----------
        x : torch.Tensor
            Input at time t, shape (batch, input_size)
        h_prev : torch.Tensor
            Previous hidden state, shape (batch, hidden_size)
 
        Returns:
        --------
        h_new : torch.Tensor
            New hidden state, shape (batch, hidden_size)
        """
        # h_t = tanh(W_ih · x_t + W_hh · h_{t-1})
        h_new = torch.tanh(self.W_ih(x) + self.W_hh(h_prev))
 
        return h_new
 
 
# Test
print("Testing VanillaRNNCell...")
cell = VanillaRNNCell(input_size=32, hidden_size=16)
x = torch.randn(4, 32)  # batch=4, input=32
h = torch.zeros(4, 16)  # initial hidden state
h_new = cell(x, h)
print(f"  Input: {x.shape}")
print(f"  Hidden: {h.shape} → {h_new.shape}")
print("✓ VanillaRNNCell works!")
\end{lstlisting}
 
\subsubsection*{6.2 Full Vanilla RNN Layer}
 
\begin{lstlisting}
class VanillaRNN(nn.Module):
    """
    Complete Vanilla RNN Layer.
 
    Επεξεργάζεται ολόκληρη την ακολουθία και επιστρέφει
    τα hidden states για κάθε timestep.
    """
 
    def __init__(self, input_size, hidden_size, num_layers=1,
                 bidirectional=False, dropout=0.0):
        super().__init__()
 
        self.input_size = input_size
        self.hidden_size = hidden_size
        self.num_layers = num_layers
        self.bidirectional = bidirectional
        self.num_directions = 2 if bidirectional else 1
 
        # RNN Cells for each layer and direction
        self.cells_forward = nn.ModuleList()
        self.cells_backward = nn.ModuleList()
 
        for layer in range(num_layers):
            layer_input_size = input_size if layer == 0 else hidden_size * self.num_directions
 
            self.cells_forward.append(VanillaRNNCell(layer_input_size, hidden_size))
 
            if bidirectional:
                self.cells_backward.append(VanillaRNNCell(layer_input_size, hidden_size))
 
        self.dropout = nn.Dropout(dropout) if dropout > 0 else None
 
    def forward(self, x, h_0=None):
        """
        Forward pass.
 
        Parameters:
        -----------
        x : torch.Tensor
            Input sequence, shape (batch, seq_len, input_size)
        h_0 : torch.Tensor, optional
            Initial hidden state
 
        Returns:
        --------
        output : torch.Tensor
            Hidden states for all timesteps, shape (batch, seq_len, hidden_size * num_directions)
        h_n : torch.Tensor
            Final hidden state
        """
        batch_size, seq_len, _ = x.shape
 
        # Initialize hidden state
        if h_0 is None:
            h_0 = torch.zeros(self.num_layers * self.num_directions,
                             batch_size, self.hidden_size, device=x.device)
 
        # Process each layer
        current_input = x
        final_hidden = []
 
        for layer in range(self.num_layers):
            # Forward direction
            h_forward = h_0[layer * self.num_directions]
            outputs_forward = []
 
            for t in range(seq_len):
                h_forward = self.cells_forward[layer](current_input[:, t, :], h_forward)
                outputs_forward.append(h_forward)
 
            outputs_forward = torch.stack(outputs_forward, dim=1)  # (batch, seq_len, hidden)
            final_hidden.append(h_forward)
 
            # Backward direction
            if self.bidirectional:
                h_backward = h_0[layer * self.num_directions + 1]
                outputs_backward = []
 
                for t in range(seq_len - 1, -1, -1):
                    h_backward = self.cells_backward[layer](current_input[:, t, :], h_backward)
                    outputs_backward.append(h_backward)
 
                outputs_backward = torch.stack(outputs_backward[::-1], dim=1)
                final_hidden.append(h_backward)
 
                # Concatenate forward and backward
                layer_output = torch.cat([outputs_forward, outputs_backward], dim=2)
            else:
                layer_output = outputs_forward
 
            # Apply dropout between layers
            if self.dropout is not None and layer < self.num_layers - 1:
                layer_output = self.dropout(layer_output)
 
            current_input = layer_output
 
        output = current_input
        h_n = torch.stack(final_hidden, dim=0)
 
        return output, h_n
 
 
# Test
print("\nTesting VanillaRNN...")
rnn = VanillaRNN(input_size=32, hidden_size=16, bidirectional=True)
x = torch.randn(4, 10, 32)  # batch=4, seq_len=10, input=32
output, h_n = rnn(x)
print(f"  Input: {x.shape}")
print(f"  Output: {output.shape}")
print(f"  Final hidden: {h_n.shape}")
print("✓ VanillaRNN works!")
\end{lstlisting}
 
\subsubsection*{6.3 Attention Mechanism}
 
\begin{lstlisting}
class Attention(nn.Module):
    """
    Attention Mechanism.
 
    Υπολογίζει ποια positions είναι σημαντικά και
    δημιουργεί weighted sum.
 
    Πλεονέκτημα: Interpretability!
    Μπορούμε να δούμε ποιες πρωτεΐνες "κοιτάει" το model.
    """
 
    def __init__(self, input_dim, attention_dim):
        super().__init__()
 
        self.attention = nn.Sequential(
            nn.Linear(input_dim, attention_dim),
            nn.Tanh(),
            nn.Linear(attention_dim, 1)
        )
 
    def forward(self, x):
        """
        Parameters:
        -----------
        x : torch.Tensor
            Shape (batch, seq_len, input_dim)
 
        Returns:
        --------
        context : torch.Tensor
            Weighted sum, shape (batch, input_dim)
        weights : torch.Tensor
            Attention weights, shape (batch, seq_len)
        """
        # Calculate attention scores
        scores = self.attention(x)  # (batch, seq_len, 1)
 
        # Softmax to get weights
        weights = F.softmax(scores, dim=1)  # (batch, seq_len, 1)
 
        # Weighted sum
        context = (x * weights).sum(dim=1)  # (batch, input_dim)
 
        return context, weights.squeeze(-1)
 
 
# Test
print("\nTesting Attention...")
attn = Attention(input_dim=32, attention_dim=16)
x = torch.randn(4, 10, 32)
context, weights = attn(x)
print(f"  Input: {x.shape}")
print(f"  Context: {context.shape}")
print(f"  Weights: {weights.shape}")
print(f"  Weights sum: {weights[0].sum().item():.4f} (should be 1.0)")
print("✓ Attention works!")
\end{lstlisting}
 
\subsubsection*{6.4 Complete Model: CNN + Vanilla RNN + Attention}
 
\begin{lstlisting}
class TemporalProteinRNN(nn.Module):
    """
    Complete Model: 1D-CNN → Vanilla RNN → Attention → Classifier
 
    Architecture:
    1. CNN: Εξάγει local patterns από τις πρωτεΐνες
    2. RNN: Μοντελοποιεί την χρονική εξέλιξη
    3. Attention: Εστιάζει στα σημαντικά features
    4. Classifier: Τελική πρόβλεψη
    """
 
    def __init__(
        self,
        n_features,
        n_timepoints=2,
        cnn_filters=32,
        cnn_kernel_size=5,
        rnn_hidden_size=32,
        rnn_num_layers=1,
        rnn_bidirectional=True,
        attention_dim=32,
        n_classes=3,
        dropout=0.5,
        use_pytorch_rnn=False  # True για PyTorch built-in, False για custom
    ):
        super().__init__()
 
        self.n_features = n_features
        self.n_timepoints = n_timepoints
        self.use_pytorch_rnn = use_pytorch_rnn
 
        # ============ 1. CNN Feature Extractor ============
        self.cnn = nn.Sequential(
            nn.Conv1d(
                in_channels=n_timepoints,
                out_channels=cnn_filters,
                kernel_size=cnn_kernel_size,
                padding=cnn_kernel_size // 2
            ),
            nn.BatchNorm1d(cnn_filters),
            nn.ReLU(),
            nn.Dropout(dropout)
        )
 
        # ============ 2. Vanilla RNN ============
        if use_pytorch_rnn:
            # PyTorch built-in RNN
            self.rnn = nn.RNN(
                input_size=cnn_filters,
                hidden_size=rnn_hidden_size,
                num_layers=rnn_num_layers,
                batch_first=True,
                bidirectional=rnn_bidirectional,
                dropout=dropout if rnn_num_layers > 1 else 0
            )
        else:
            # Custom Vanilla RNN
            self.rnn = VanillaRNN(
                input_size=cnn_filters,
                hidden_size=rnn_hidden_size,
                num_layers=rnn_num_layers,
                bidirectional=rnn_bidirectional,
                dropout=dropout if rnn_num_layers > 1 else 0
            )
 
        # RNN output size
        rnn_output_size = rnn_hidden_size * (2 if rnn_bidirectional else 1)
 
        # ============ 3. Attention ============
        self.attention = Attention(rnn_output_size, attention_dim)
 
        # ============ 4. Classifier ============
        self.classifier = nn.Sequential(
            nn.Dropout(dropout),
            nn.Linear(rnn_output_size, n_classes)
        )
 
        # Store attention weights for interpretability
        self.attention_weights = None
 
    def forward(self, x):
        """
        Forward pass.
 
        Parameters:
        -----------
        x : torch.Tensor
            Shape (batch, n_timepoints, n_features)
 
        Returns:
        --------
        logits : torch.Tensor
            Shape (batch, n_classes)
        """
        # x: (batch, timepoints, features)
 
        # 1. CNN: Extract local patterns
        # CNN expects (batch, channels, length)
        cnn_out = self.cnn(x)  # (batch, cnn_filters, seq_len)
 
        # 2. RNN: Process sequence
        # RNN expects (batch, seq_len, features)
        rnn_in = cnn_out.transpose(1, 2)  # (batch, seq_len, cnn_filters)
        rnn_out, _ = self.rnn(rnn_in)  # (batch, seq_len, rnn_hidden * directions)
 
        # 3. Attention: Focus on important parts
        context, attn_weights = self.attention(rnn_out)  # (batch, rnn_hidden * directions)
        self.attention_weights = attn_weights  # Store for interpretability
 
        # 4. Classifier: Final prediction
        logits = self.classifier(context)  # (batch, n_classes)
 
        return logits
 
    def get_attention_weights(self):
        """Returns the last computed attention weights"""
        return self.attention_weights
\end{lstlisting}
 
\begin{lstlisting}
# Test the complete model
print("Testing Complete Model...")
print("="*50)
 
n_features = X.shape[2]
n_classes = len(class_names)
 
# Create model
model = TemporalProteinRNN(
    n_features=n_features,
    n_timepoints=2,
    cnn_filters=Config.CNN_FILTERS,
    cnn_kernel_size=Config.CNN_KERNEL_SIZE,
    rnn_hidden_size=Config.RNN_HIDDEN_SIZE,
    rnn_num_layers=Config.RNN_NUM_LAYERS,
    rnn_bidirectional=Config.RNN_BIDIRECTIONAL,
    attention_dim=Config.ATTENTION_DIM,
    n_classes=n_classes,
    dropout=Config.DROPOUT,
    use_pytorch_rnn=False  # Use our custom Vanilla RNN
)
 
# Count parameters
total_params = sum(p.numel() for p in model.parameters())
trainable_params = sum(p.numel() for p in model.parameters() if p.requires_grad)
 
print(f"\nModel Architecture:")
print(f"  Input: ({Config.BATCH_SIZE}, 2, {n_features})")
print(f"  CNN filters: {Config.CNN_FILTERS}")
print(f"  RNN hidden: {Config.RNN_HIDDEN_SIZE}")
print(f"  Bidirectional: {Config.RNN_BIDIRECTIONAL}")
print(f"  Output: ({Config.BATCH_SIZE}, {n_classes})")
print(f"\n  Total parameters: {total_params:,}")
print(f"  Trainable parameters: {trainable_params:,}")
print(f"  Samples / Parameters ratio: {len(y) / total_params:.3f}")
 
# Test forward pass
dummy_input = torch.randn(4, 2, n_features)
dummy_output = model(dummy_input)
print(f"\nForward pass test:")
print(f"  Input shape: {dummy_input.shape}")
print(f"  Output shape: {dummy_output.shape}")
print(f"  Attention weights shape: {model.get_attention_weights().shape}")
 
print("\n✓ Model ready!")
\end{lstlisting}
 
\subsubsection*{ΜΕΡΟΣ 7: TRAINING FUNCTIONS}
 
\begin{lstlisting}
def train_epoch(model, loader, optimizer, criterion, device):
    """
    Training για ένα epoch.
 
    Returns:
    --------
    avg_loss : float
    accuracy : float
    """
    model.train()
    total_loss = 0
    correct = 0
    total = 0
 
    for X_batch, y_batch in loader:
        X_batch = X_batch.to(device)
        y_batch = y_batch.to(device)
 
        # Forward pass
        optimizer.zero_grad()
        outputs = model(X_batch)
        loss = criterion(outputs, y_batch)
 
        # Check for NaN
        if torch.isnan(loss):
            print("WARNING: NaN loss!")
            continue
 
        # Backward pass
        loss.backward()
 
        # Gradient clipping (prevents exploding gradients)
        torch.nn.utils.clip_grad_norm_(model.parameters(), max_norm=1.0)
 
        # Update weights
        optimizer.step()
 
        # Statistics
        total_loss += loss.item()
        _, predicted = outputs.max(1)
        correct += (predicted == y_batch).sum().item()
        total += y_batch.size(0)
 
    avg_loss = total_loss / len(loader) if len(loader) > 0 else 0
    accuracy = correct / total if total > 0 else 0
 
    return avg_loss, accuracy
\end{lstlisting}
 
\begin{lstlisting}
def evaluate(model, loader, criterion, device):
    """
    Evaluation.
 
    Returns:
    --------
    avg_loss : float
    accuracy : float
    predictions : np.ndarray
    targets : np.ndarray
    """
    model.eval()
    total_loss = 0
    all_preds = []
    all_targets = []
 
    with torch.no_grad():
        for X_batch, y_batch in loader:
            X_batch = X_batch.to(device)
            y_batch = y_batch.to(device)
 
            # Forward pass
            outputs = model(X_batch)
            loss = criterion(outputs, y_batch)
 
            if not torch.isnan(loss):
                total_loss += loss.item()
 
            # Predictions
            _, predicted = outputs.max(1)
            all_preds.extend(predicted.cpu().numpy())
            all_targets.extend(y_batch.cpu().numpy())
 
    avg_loss = total_loss / len(loader) if len(loader) > 0 else 0
    accuracy = accuracy_score(all_targets, all_preds)
 
    return avg_loss, accuracy, np.array(all_preds), np.array(all_targets)
 
print("✓ Training functions defined!")
\end{lstlisting}
 
\subsubsection*{ΜΕΡΟΣ 8: CROSS-VALIDATION}
 
\begin{lstlisting}
def run_cross_validation(X, y, mode, config, n_folds=5, use_pytorch_rnn=False):
    """
    Stratified K-Fold Cross-Validation.
 
    Parameters:
    -----------
    X : np.ndarray
        Data
    y : np.ndarray
        Labels
    mode : str
        'temporal', 'static', or 'delta'
    config : Config
        Configuration object
    n_folds : int
        Number of CV folds
    use_pytorch_rnn : bool
        Use PyTorch RNN or custom implementation
 
    Returns:
    --------
    results : dict
        Results including accuracy, predictions, etc.
    """
    print(f"\n{'='*60}")
    print(f"EXPERIMENT: {mode.upper()} MODEL")
    print(f"{'='*60}")
    print(f"Mode: {mode}")
    print(f"Folds: {n_folds}")
    print(f"RNN: {'PyTorch' if use_pytorch_rnn else 'Custom Vanilla'}")
 
    # Determine n_timepoints based on mode
    n_timepoints = 2 if mode == 'temporal' else 1
    n_features = X.shape[2]
    n_classes = len(np.unique(y))
 
    # Class weights for imbalanced data
    class_counts = np.bincount(y)
    class_weights = torch.FloatTensor(
        len(y) / (len(class_counts) * class_counts)
    ).to(config.DEVICE)
 
    # Cross-validation
    kfold = StratifiedKFold(n_splits=n_folds, shuffle=True, random_state=config.RANDOM_SEED)
 
    fold_results = []
    all_preds = []
    all_targets = []
    histories = []
 
    for fold, (train_idx, val_idx) in enumerate(kfold.split(X, y)):
        print(f"\n--- Fold {fold+1}/{n_folds} ---")
 
        # Split data
        X_train, X_val = X[train_idx], X[val_idx]
        y_train, y_val = y[train_idx], y[val_idx]
 
        # Create datasets
        train_dataset = ProteinDataset(X_train, y_train, mode=mode)
        val_dataset = ProteinDataset(X_val, y_val, mode=mode)
 
        train_loader = DataLoader(
            train_dataset,
            batch_size=config.BATCH_SIZE,
            shuffle=True
        )
        val_loader = DataLoader(
            val_dataset,
            batch_size=config.BATCH_SIZE
        )
 
        # Create model
        model = TemporalProteinRNN(
            n_features=n_features,
            n_timepoints=n_timepoints,
            cnn_filters=config.CNN_FILTERS,
            cnn_kernel_size=config.CNN_KERNEL_SIZE,
            rnn_hidden_size=config.RNN_HIDDEN_SIZE,
            rnn_num_layers=config.RNN_NUM_LAYERS,
            rnn_bidirectional=config.RNN_BIDIRECTIONAL,
            attention_dim=config.ATTENTION_DIM,
            n_classes=n_classes,
            dropout=config.DROPOUT,
            use_pytorch_rnn=use_pytorch_rnn
        ).to(config.DEVICE)
 
        # Loss and optimizer
        criterion = nn.CrossEntropyLoss(weight=class_weights)
        optimizer = torch.optim.AdamW(
            model.parameters(),
            lr=config.LEARNING_RATE,
            weight_decay=config.WEIGHT_DECAY
        )
        scheduler = torch.optim.lr_scheduler.ReduceLROnPlateau(
            optimizer, mode='min', factor=0.5, patience=5
        )
 
        # Training loop
        best_val_acc = 0
        best_model_state = None
        patience_counter = 0
        history = {'train_loss': [], 'train_acc': [], 'val_loss': [], 'val_acc': []}
 
        for epoch in range(config.EPOCHS):
            # Train
            train_loss, train_acc = train_epoch(
                model, train_loader, optimizer, criterion, config.DEVICE
            )
 
            # Validate
            val_loss, val_acc, _, _ = evaluate(
                model, val_loader, criterion, config.DEVICE
            )
 
            # Update scheduler
            scheduler.step(val_loss)
 
            # Save history
            history['train_loss'].append(train_loss)
            history['train_acc'].append(train_acc)
            history['val_loss'].append(val_loss)
            history['val_acc'].append(val_acc)
 
            # Check for improvement
            if val_acc > best_val_acc:
                best_val_acc = val_acc
                best_model_state = {k: v.cpu().clone() for k, v in model.state_dict().items()}
                patience_counter = 0
            else:
                patience_counter += 1
 
            # Early stopping
            if patience_counter >= config.PATIENCE:
                print(f"  Early stopping at epoch {epoch+1}")
                break
 
            # Print progress
            if (epoch + 1) % 20 == 0:
                print(f"  Epoch {epoch+1}: Train Acc={train_acc:.4f}, Val Acc={val_acc:.4f}")
 
        # Load best model and evaluate
        if best_model_state is not None:
            model.load_state_dict(best_model_state)
 
        _, val_acc, preds, targets = evaluate(
            model, val_loader, criterion, config.DEVICE
        )
 
        fold_results.append(val_acc)
        all_preds.extend(preds)
        all_targets.extend(targets)
        histories.append(history)
 
        print(f"  Fold {fold+1} Best Accuracy: {val_acc:.4f}")
 
    # Summary
    mean_acc = np.mean(fold_results)
    std_acc = np.std(fold_results)
 
    print(f"\n{'='*40}")
    print(f"RESULTS: {mode.upper()}")
    print(f"{'='*40}")
    print(f"Mean Accuracy: {mean_acc:.4f} ± {std_acc:.4f}")
    print(f"Per-fold: {[f'{acc:.3f}' for acc in fold_results]}")
 
    return {
        'mean_acc': mean_acc,
        'std_acc': std_acc,
        'fold_results': fold_results,
        'all_preds': np.array(all_preds),
        'all_targets': np.array(all_targets),
        'histories': histories
    }
\end{lstlisting}
 
\subsubsection*{ΜΕΡΟΣ 9: RUN EXPERIMENTS ��}
 
\begin{lstlisting}
# Run all experiments
print("\n" + "#"*60)
print("#" + " "*20 + "RUNNING EXPERIMENTS" + " "*19 + "#")
print("#"*60)
 
results = {}
 
for mode in ['static', 'delta', 'temporal']:
    results[mode] = run_cross_validation(
        X, y,
        mode=mode,
        config=Config,
        n_folds=Config.N_FOLDS,
        use_pytorch_rnn=False  # Use custom Vanilla RNN
    )
\end{lstlisting}
 
\subsubsection*{ΜΕΡΟΣ 10: RESULTS \& VISUALIZATION}
 
\begin{lstlisting}
# Summary Table
print("\n" + "="*60)
print("FINAL RESULTS COMPARISON")
print("="*60)
print(f"\n{'Model':<15} {'Mean Accuracy':<18} {'Std':<10} {'Per Fold'}")
print("-" * 70)
 
for mode in ['static', 'delta', 'temporal']:
    mean = results[mode]['mean_acc']
    std = results[mode]['std_acc']
    folds = [f"{x:.3f}" for x in results[mode]['fold_results']]
    print(f"{mode:<15} {mean:<18.4f} ±{std:<9.4f} {folds}")
 
# Random baseline
random_baseline = 1.0 / len(class_names)
print(f"\n{'random guess':<15} {random_baseline:<18.4f}")
 
# Conclusion
print("\n" + "="*60)
print("CONCLUSION")
print("="*60)
 
temporal_acc = results['temporal']['mean_acc']
static_acc = results['static']['mean_acc']
delta_acc = results['delta']['mean_acc']
 
best_mode = max(['static', 'delta', 'temporal'], key=lambda m: results[m]['mean_acc'])
print(f"\nBest Model: {best_mode.upper()} ({results[best_mode]['mean_acc']:.4f})")
 
if temporal_acc > static_acc + 0.02:
    diff = (temporal_acc - static_acc) * 100
    print(f"\n✓ HYPOTHESIS SUPPORTED!")
    print(f"  Temporal model outperforms Static by {diff:.1f}%")
    print(f"  → Temporal information adds diagnostic value!")
elif static_acc > temporal_acc + 0.02:
    print(f"\n⚠ Static model performs better")
    print(f"  → May need more data for temporal patterns")
else:
    print(f"\n~ Results are similar (within noise)")
    print(f"  → Need more data for definitive conclusions")
\end{lstlisting}
 
\begin{lstlisting}
# Visualization
fig, axes = plt.subplots(1, 2, figsize=(14, 5))
 
# 1. Bar plot comparison
modes = ['static', 'delta', 'temporal']
means = [results[m]['mean_acc'] for m in modes]
stds = [results[m]['std_acc'] for m in modes]
 
colors = ['#3498db', '#e74c3c', '#2ecc71']
x_pos = np.arange(len(modes))
 
bars = axes[0].bar(x_pos, means, yerr=stds, capsize=5, color=colors, alpha=0.8, edgecolor='black')
axes[0].set_xticks(x_pos)
axes[0].set_xticklabels(['Static\n(T0 only)', 'Delta\n(T1-T0)', 'Temporal\n([T0,T1])'])
axes[0].set_ylabel('Accuracy', fontsize=12)
axes[0].set_title('Vanilla RNN: Model Comparison\n(5-Fold Cross-Validation)', fontsize=14)
axes[0].set_ylim(0, 1)
axes[0].axhline(y=1/len(class_names), color='gray', linestyle='--', label='Random guess', alpha=0.7)
axes[0].legend(loc='upper right')
axes[0].grid(axis='y', alpha=0.3)
 
# Add values on bars
for bar, mean, std in zip(bars, means, stds):
    axes[0].text(bar.get_x() + bar.get_width()/2, mean + std + 0.03,
                 f'{mean:.3f}', ha='center', fontsize=11, fontweight='bold')
 
# 2. Confusion matrix for best model
best_mode = max(modes, key=lambda m: results[m]['mean_acc'])
cm = confusion_matrix(results[best_mode]['all_targets'], results[best_mode]['all_preds'])
 
sns.heatmap(cm, annot=True, fmt='d', cmap='Blues',
            xticklabels=class_names, yticklabels=class_names, ax=axes[1],
            annot_kws={"size": 14})
axes[1].set_xlabel('Predicted', fontsize=12)
axes[1].set_ylabel('True', fontsize=12)
axes[1].set_title(f'Confusion Matrix\n({best_mode.capitalize()} Model)', fontsize=14)
 
plt.tight_layout()
plt.show()
\end{lstlisting}
 
\begin{lstlisting}
# Classification Report
best_mode = max(modes, key=lambda m: results[m]['mean_acc'])
 
print(f"\nClassification Report ({best_mode.capitalize()} Model):")
print("="*55)
print(classification_report(
    results[best_mode]['all_targets'],
    results[best_mode]['all_preds'],
    target_names=class_names
))
\end{lstlisting}
 
\begin{lstlisting}
# Training History (from last fold)
fig, axes = plt.subplots(1, 3, figsize=(15, 4))
 
for idx, mode in enumerate(['static', 'delta', 'temporal']):
    history = results[mode]['histories'][-1]  # Last fold
 
    axes[idx].plot(history['train_acc'], label='Train', linewidth=2)
    axes[idx].plot(history['val_acc'], label='Validation', linewidth=2)
    axes[idx].set_xlabel('Epoch')
    axes[idx].set_ylabel('Accuracy')
    axes[idx].set_title(f'{mode.capitalize()} Model')
    axes[idx].legend()
    axes[idx].grid(True, alpha=0.3)
    axes[idx].set_ylim(0, 1)
 
plt.suptitle('Training History (Last Fold)', fontsize=14, y=1.02)
plt.tight_layout()
plt.show()
\end{lstlisting}
 
\subsubsection*{ΜΕΡΟΣ 11: TOP DISCRIMINATIVE PROTEINS}
 
\begin{lstlisting}
# Get protein importance from feature selection
scores = feature_selector.scores_[selected_indices]
protein_importance = sorted(
    zip(selected_proteins, scores),
    key=lambda x: x[1],
    reverse=True
)
 
print("\nTop 20 Most Discriminative Proteins:")
print("="*50)
for i, (protein, score) in enumerate(protein_importance[:20]):
    print(f"{i+1:2d}. {protein:<20} (F-score: {score:.2f})")
\end{lstlisting}
 
\begin{lstlisting}
# Visualize
top_n = 15
top_proteins = [p for p, _ in protein_importance[:top_n]]
top_scores = [s for _, s in protein_importance[:top_n]]
 
fig, ax = plt.subplots(figsize=(10, 6))
bars = ax.barh(range(top_n), top_scores[::-1], color='steelblue', alpha=0.8, edgecolor='black')
ax.set_yticks(range(top_n))
ax.set_yticklabels(top_proteins[::-1])
ax.set_xlabel('F-score (ANOVA)', fontsize=12)
ax.set_title('Top 15 Discriminative Proteins', fontsize=14)
ax.grid(axis='x', alpha=0.3)
 
plt.tight_layout()
plt.show()
\end{lstlisting}
 
\subsubsection*{ΜΕΡΟΣ 12: SAVE RESULTS}
 
\begin{lstlisting}
# Save results to CSV
results_df = pd.DataFrame({
    'Model': ['Static', 'Delta', 'Temporal'],
    'Mean_Accuracy': [results['static']['mean_acc'], results['delta']['mean_acc'], results['temporal']['mean_acc']],
    'Std_Accuracy': [results['static']['std_acc'], results['delta']['std_acc'], results['temporal']['std_acc']],
})
 
# Add fold results
for i in range(Config.N_FOLDS):
    results_df[f'Fold_{i+1}'] = [
        results['static']['fold_results'][i],
        results['delta']['fold_results'][i],
        results['temporal']['fold_results'][i]
    ]
 
results_df.to_csv('vanilla_rnn_results.csv', index=False)
print("Results saved to 'vanilla_rnn_results.csv'")
display(results_df)
\end{lstlisting}
 
\begin{lstlisting}
# Save top proteins
proteins_df = pd.DataFrame({
    'Rank': range(1, len(protein_importance) + 1),
    'Protein': [p for p, _ in protein_importance],
    'F_score': [s for _, s in protein_importance]
})
 
proteins_df.to_csv('top_proteins_vanilla_rnn.csv', index=False)
print("\nTop proteins saved to 'top_proteins_vanilla_rnn.csv'")
display(proteins_df.head(15))
\end{lstlisting}
 
\subsubsection*{SUMMARY}
 
